%! TEX root = main.tex  % Specifies the main document file for multi-file projects

%%%%%%%%%%%%%%%%%%%%%%%%%%%%%%%%%%%%%%%%%%%%%%%%%%%%%%%%%%%%%%%%%%%%%%%%%%%%%%%
% Custom Colors, Math Symbols, and Highlighting Command
%%%%%%%%%%%%%%%%%%%%%%%%%%%%%%%%%%%%%%%%%%%%%%%%%%%%%%%%%%%%%%%%%%%%%%%%%%%%%%%
% Define custom colors for use throughout the document.
\definecolor{bluePoliMI}{cmyk}{0.4, 0.1, 0, 0.4}
\definecolor{redUniPD}{RGB}{155, 0, 20}
\definecolor{pythonGreen}{RGB}{0, 127, 0}
\definecolor{specialGreen}{RGB}{50, 255, 120}
\definecolor{bluePoliTO}{RGB}{0,86,145}

% Define a command to highlight text with a light gray background.
\newcommand{\hl}[1]{\colorbox{gray!10}{#1}}

% Explicitly declare the plus sign as a math binary operator.
%\DeclareMathSymbol{+}{\mathbin}{operators}{"2B}

%%%%%%%%%%%%%%%%%%%%%%%%%%%%%%%%%%%%%%%%%%%%%%%%%%%%%%%%%%%%%%%%%%%%%%%%%%%%%%%
% Theorem Environments Setup
%%%%%%%%%%%%%%%%%%%%%%%%%%%%%%%%%%%%%%%%%%%%%%%%%%%%%%%%%%%%%%%%%%%%%%%%%%%%%%%
% Define custom theorem style with blue titles and italic body text
\declaretheoremstyle[
    headfont=\color{bluePoliMI}\normalfont\bfseries,
    bodyfont=\color{black}\normalfont\itshape,
]{colored}
\theoremstyle{colored}
\newtheorem{theorem}{Theorem}[section]       % Theorem environment, resets each section
\newtheorem{proposition}{Proposition}[section] % Proposition environment, resets each section

%%%%%%%%%%%%%%%%%%%%%%%%%%%%%%%%%%%%%%%%%%%%%%%%%%%%%%%%%%%%%%%%%%%%%%%%%%%%%%%
% Table Column Types and Spacing Commands
%%%%%%%%%%%%%%%%%%%%%%%%%%%%%%%%%%%%%%%%%%%%%%%%%%%%%%%%%%%%%%%%%%%%%%%%%%%%%%%
% Define custom column types:
% - L: Left-aligned, fixed width, allows breaks
% - C: Centered, fixed width
% - R: Right-aligned, fixed width
\newcolumntype{L}[1]{>{\raggedright\let\newline\\\arraybackslash\hspace{0pt}}m{#1}}
\newcolumntype{C}[1]{>{\centering\let\newline\\\arraybackslash\hspace{0pt}}m{#1}}
\newcolumntype{R}[1]{>{\raggedleft\let\newline\\\arraybackslash\hspace{0pt}}m{#1}}

% Commands for vertical spacing in table cells
\newcommand\T{\rule{0pt}{2.6ex}}
\newcommand\B{\rule[-1.2ex]{0pt}{0pt}}

%%%%%%%%%%%%%%%%%%%%%%%%%%%%%%%%%%%%%%%%%%%%%%%%%%%%%%%%%%%%%%%%%%%%%%%%%%%%%%%
% List Formatting Settings
%%%%%%%%%%%%%%%%%%%%%%%%%%%%%%%%%%%%%%%%%%%%%%%%%%%%%%%%%%%%%%%%%%%%%%%%%%%%%%%
% Customize itemize lists:
% - Level 1: Bullet
% - Level 2: Circle
% - Level 3: Dash
% Remove extra item spacing
\setlist[itemize,1]{label=\( \bullet \)}
\setlist[itemize,2]{label=\( \circ \)}
\setlist[itemize,3]{label=\( - \)}
\setlist{nosep}

%%%%%%%%%%%%%%%%%%%%%%%%%%%%%%%%%%%%%%%%%%%%%%%%%%%%%%%%%%%%%%%%%%%%%%%%%%%%%%%
% General Layout Settings
%%%%%%%%%%%%%%%%%%%%%%%%%%%%%%%%%%%%%%%%%%%%%%%%%%%%%%%%%%%%%%%%%%%%%%%%%%%%%%%
\setlength\parindent{0pt}  % Remove paragraph indentation

% Custom chapter heading inside a colored box
\newcommand{\chapterbox}[1]{
    \tikz[baseline=(box.base)]{
        \node[rectangle, rounded corners=2pt, fill=bluePoliMI, text=white, inner xsep=10pt, inner ysep=7pt](box){#1};
    }
}

% Chapter title formatting
\titleformat{\chapter}[display]
    {\huge\bfseries}            % Title style
    {\chapterbox{\thechapter}}  % Chapter number inside a box
    {1ex}                        % Vertical space
    {\scshape}                  % Text style

% Custom section heading inside a colored box
\newcommand{\sectionbox}[1]{
    \tikz[baseline=(box.base)]{
        \node[rectangle, rounded corners=1pt, fill=bluePoliMI, text=white, inner xsep=4pt, inner ysep=3pt](box){#1};
    } 
}

% Section title formatting
\titleformat{\section}[hang]
    {\Large\bfseries}          % Title style
    {\sectionbox{\thesection}} % Section number inside a box
    {1em}                       % Horizontal space
    {\scshape}                 % Text style

% Uncomment to order section with letter
%\renewcommand{\thesection}{\thechapter.\alph{section}}

\titleformat{\subsection}[hang]
    {\large\bfseries}          % Subsection title style
    {\sectionbox{\thesubsection}} % Reuse section box
    {1em}                       % Space between number and text
    {\scshape}                 % Text style

% Uncomment to order subsection romanically
%\renewcommand{\thesubsection}{\thesection.\roman{subsection}}

\titleformat{\subsubsection}[hang]
    {\normalsize}              % Subsubsection title style
    {}{0ex}{\bfseries\scshape} % No number box, styled text

% Customization of the \@chapter command to include numbered chapters in TOC and LoL
\makeatletter
\def\@chapter[#1]#2{
    \ifnum \c@secnumdepth >\m@ne
        \if@mainmatter
            \refstepcounter{chapter} % Increment chapter counter
            \typeout{\@chapapp\space\thechapter.} % Print chapter info to log
            \addcontentsline{toc}{chapter}{\protect\numberline{\thechapter}#1} % Add to TOC
        \else
            \addcontentsline{toc}{chapter}{#1} % Unnumbered chapter in TOC
        \fi
    \else
        \addcontentsline{toc}{chapter}{#1} % Unnumbered chapter in TOC
    \fi
    \chaptermark{#1} % Set chapter header
    \addtocontents{lof}{\protect\addvspace{10\p@}} % Add space in LoF
    \addtocontents{lot}{\protect\addvspace{10\p@}} % Add space in LoT
    \if@twocolumn
        \@topnewpage[\@makechapterhead{#2}] % Format for two-column layout
    \else
        \@makechapterhead{#2} % Standard chapter heading
        \@afterheading
    \fi
}

% Patch \@chapter to add spacing for chapter entries in the List of Listings (LoL)
\patchcmd{\@chapter}% <cmd>
  {\addtocontents}% <search>
  {\addtocontents{lol}{\protect\addvspace{10\p@}}% Add per-chapter space in LoL
   \addtocontents}% <replace>
  {}{}% <success><failure>
\makeatother